\section{Introduction}

In this paper, we study rigidity results for metrics with lower scalar curvature bounds.
One of the first results of this kind is the famous rigidity theorem of Llarull \cite{Llarull}.
Let $g_{S^n}$ denote the standard round metric on $S^n$ with scalar curvature $n(n-1)$. Llarull showed that, if $g$ is a metric on $S^n$ with $g \geq g_{S^n}$ and $R_g \geq n(n-1)$, then $g=g_{S^n}$.
The proof of Llarull's theorem uses Dirac operator techniques in an ingenious way, and is inspired by the fundamental work of Gromov and Lawson \cite{Gromov-Lawson1, Gromov-Lawson2}.

In an important paper, Cecchini and Zeidler extended this line of thought and proved scalar and mean curvature rigidity results for odd-dimensional manifolds with boundary where the comparison metric is not the round metric, but a warped product metric (see \cite[Section~10]{Cecchini-Zeidler}).
In the first part of the present paper we will remove the dimension parity assumption in some of their results.

Given a Riemannian manifold $(M,g)$ with boundary, we denote by $R_g$ the scalar curvature of $g$. Moreover, we denote by $\nu_g$ the outward unit normal with respect to $g$.
We denote by $H_g$ the mean curvature of $\partial M$ with respect to $g$, defined as the sum of the principal curvatures.
The sign convention for $H_g$ is such that the mean curvature vector is given by $-H_g\nu_g$.

Let $n>2$, let $\theta_- < \theta_+$ and let $\rho\colon [\theta_- , \theta_+] \to \R$ be a positive smooth function.
We consider the warped product metric
\begin{equation} \label{g_0}
 g_0 = {\rm d}\theta \otimes {\rm d}\theta + \rho^2(\theta) g_{S^{n-1}}
\end{equation}
on $S^{n-1} \times [\theta_- , \theta_+]$. The scalar curvature of $g_0$ is given by
\begin{equation} \label{scal_g_0}
 R_{g_0} = (n-1) \left( -2 \frac{\rho''(\theta)}{\rho(\theta)} + (n-2) \frac{1- \rho'(\theta)^2}{\rho(\theta)^2} \right),
\end{equation}
while the boundary mean curvature of $g_0$ is given by
\begin{equation} \label{mean_g_0}
 H_{g_0} = \pm (n-1) \frac{\rho'(\theta_{\pm})}{\rho(\theta_{\pm})}\qquad \text{along}\quad S^{n-1} \times \{ \theta_{\pm} \}
\end{equation}
(cf.\ \cite[Example~4.1]{Baer-Gauduchon-Moroianu}).

Our first result says that warped product metrics satisfy a scalar-mean curvature rigidity property, provided that the warping function is strictly logarithmically concave.

\begin{theorem} \label{spherical_band}
Let $n>2$, let $\rho\colon [\theta_- , \theta_+] \to \R$ be a positive smooth function such that ${(\log \rho)''<0}$.
Let $g_0$ denote the warped product metric in \eqref{g_0}.
Let $M$ be a compact, connected spin manifold of dimension $n$ with boundary $\partial M$.
Let $g$ be a Riemannian metric on $M$.
Suppose that $\Phi\colon (M,g) \to \big(S^{n-1} \times [\theta_-,\theta_+],g_0\big)$ is a smooth map with the following properties:
\begin{itemize}\itemsep=0pt
\item $\Phi(\partial M) \subset S^{n-1} \times \{\theta_+,\theta_-\}$,
\item $\Phi$ has non-zero degree,
\item $\Phi$ is $1$-Lipschitz,
\item $R_g \geq R_{g_0} \circ \Phi$ at each point in $M$, compare \eqref{scal_g_0},
\item $H_g \geq H_{g_0} \circ \Phi$ at each point in $\partial M$, compare \eqref{mean_g_0}.
\end{itemize}
Then $\Phi$ is a Riemannian isometry.
\end{theorem}

We note that $R_{g_0}$ in this theorem is not required to be non-negative.
For $n$ odd, Theorem~\ref{spherical_band} is implied by results of Cecchini--Zeidler, see \cite[Theorem~10.2]{Cecchini-Zeidler}.

Applying this discussion to annuli in simply-connected space forms as in \cite[Section~10]{Cecchini-Zeidler}, this removes the parity restriction in \cite[Corollaries~10.4 and~10.5]{Cecchini-Zeidler}.

